\begin{explanationbox}

	\textbf{\textcolor{CExplanation}{一句话直觉}}

	椭圆绕  $y$ 轴旋转后形成旋转椭球。它的横向半径为 $a$，竖直半径为 $b$。当 $b>a$ 时像橄榄球，当 $a>b$ 时像扁球,它的体积等于 $a^2 b$ 乘以一个常数 $\frac{4}{3}\pi$，即体积与赤道半径的平方乘以极半径成正比。

	\medskip
	\textbf{\textcolor{CExplanation}{核心拆解}}
	\begin{description}[style=nextline, leftmargin=2.8em, labelsep=0.5em, itemsep=0.45em]
		\item[\textbf{要点一（旋转体体积公式）：}] 旋转体体积可由积分 $V=\pi\int x^2\,dy$ 计算，其中 $x$ 是截面半径。
		\item[\textbf{要点二（椭圆与半径关系）：}] 椭圆方程 $\frac{x^2}{a^2}+\frac{y^2}{b^2}=1$ 给出 $x^2=a^2\left(1-\frac{y^2}{b^2}\right)$，截面半径随 $y$ 变化。
		\item[\textbf{要点三（定积分与常数）：}] 积分 $\int_{-b}^{b}\left(1-\frac{y^2}{b^2}\right)dy$ 计算得 $\frac{4}{3}b$，乘以 $\pi a^2$ 即得 $V_y$。
	\end{description}

	\medskip
	\textbf{\textcolor{CExplanation}{几何本质（截面累加原理）}}

	旋转椭球可以看作一系列垂直于旋转轴的圆盘叠加而成。每个圆盘半径 $r(y)=|x|$，其中 $|x|=\frac{a}{b}\sqrt{b^2-y^2}$，其面积 $\pi r^2$ 在 $y$ 方向上积分即得体积。

	\medskip
	\textbf{\textcolor{CExplanation}{代数意义（公式对称性）}}

	绕哪条轴，哪条轴方向的半轴只出现一次；垂直于旋转轴的半轴要平方。绕 $y$ 轴体积正比于 $a^2b$，绕 $x$ 轴体积正比于 $ab^2$，体现出公式对半轴的对称性。

	\medskip
	\textbf{\textcolor{CExplanation}{考点价值（高考常考点）}}

	\begin{itemize}[leftmargin=*, itemsep=0.5em, topsep=0.4em]
		\item \textbf{考法一（直接计算体积）：} 给出椭圆方程，直接代入公式求旋转体体积。
		\item \textbf{考法二（反求参数）：} 利用已知体积与部分半轴，通过方程求解另一半轴。
	\end{itemize}

	\medskip
	\textbf{\textcolor{CExplanation}{顿悟点}}

	\textbf{体积由$a^2b$这个乘积决定，与$a$、$b$具体比例无关，只有当$a^2b$相同的时候，体积才相同。}

	\medskip
	\textbf{\textcolor{CExplanation}{使用场景}}
	\begin{itemize}[leftmargin=*, itemsep=0.5em, topsep=0.4em]
		\item \textbf{场景一（旋转体体积计算）：} 已知椭圆方程，求绕 $y$ 轴旋转体体积。
		\item \textbf{场景二（椭球体积比较）：} 与其它旋转体（如球、圆柱）体积比较或求比例。
	\end{itemize}


    \medskip
    \textbf{\textcolor{CExplanation}{新增视角一（换一个方向切片：沿 $x$ 轴累加）}}

    \textbf{切片方向不必等于旋转轴方向。} 原来的方法沿 $y$ 轴累加，切面 $y=y_0$ 垂直于旋转轴，所以截面是\textbf{圆}；现在改沿 $x$ 轴累加，切面 $x=x_0$ 平行于旋转轴，所以截面是\textbf{椭圆}。\textbf{千万不能把后一种截面也当成圆。}

    旋转后立体的三维方程是
    \[
      \frac{x^2}{a^2}+\frac{z^2}{a^2}+\frac{y^2}{b^2}\le 1.
    \]
    固定 $x=x_0$，设 $q=\sqrt{1-\frac{x_0^2}{a^2}}$，则截面满足
    \[
       \frac{y^2}{(bq)^2}+\frac{z^2}{(aq)^2}\le1;
       \qquad A(x_0)=\pi\,(bq)(aq)=\pi ab\left(1-\frac{x_0^2}{a^2}\right).
    \]
    因而 $V_y=\int_{-a}^{a}A(x)\,dx$。这与沿 $y$ 方向的圆盘积分是\textbf{同一立体、两种切法}。

    \begin{center}
      \resizebox{0.98\linewidth}{!}{% C026_slices.tikz.tex -- 直接在已加载 TikZ 的 LaTeX 文档中 \input 本文件。
% 结构示意：母线取 a=3, b=2（无量纲），截面示意图各自独立缩放。
\begin{tikzpicture}[x=0.72cm,y=0.72cm,>=stealth, every node/.style={font=\small}]
  % -------- 1) Meridian ellipse --------
  \begin{scope}[shift={(0,0)}]
    \fill[blue!3,rounded corners=7pt] (-3.65,-3.15) rectangle (3.65,3.05);
    \draw[blue!17,rounded corners=7pt] (-3.65,-3.15) rectangle (3.65,3.05);
    \draw[->,gray!75] (-3.35,0)--(3.35,0) node[below left,font=\scriptsize] {$x$};
    \draw[->,gray!75] (0,-2.52)--(0,2.67) node[above left,font=\scriptsize] {$y$};
    \fill[blue!10] (0,0) ellipse [x radius=3,y radius=2];
    \draw[blue!75!black,thick] (0,0) ellipse [x radius=3,y radius=2];
    % y=0.8 gives radius 3*sqrt(1-0.8^2/4)=2.74955
    \draw[orange!85!black,line width=2.6pt] (-2.7495,0.8)--(2.7495,0.8);
    \draw[orange!75!black,dashed] (-2.7495,0.8)--(-3.3,0.8);
    \node[anchor=east,text=orange!70!black,font=\scriptsize] at (-3.32,0.8) {$y_0$};
    % x=1.5 gives meridian y=+-1.732
    \draw[teal!75!black,line width=2.6pt] (1.5,-1.732)--(1.5,1.732);
    \draw[teal!75!black,dashed] (1.5,-1.732)--(1.5,-2.4);
    \node[anchor=north,text=teal!65!black,font=\scriptsize] at (1.5,-2.4) {$x_0$};
    \node[font=\scriptsize,fill=blue!3,inner sep=1pt] at (0,-2.82) {原椭圆：$\frac{x^2}{a^2}+\frac{y^2}{b^2}=1$};
    \node[text=blue!75!black,font=\scriptsize] at (2.65,0.26) {$a$};
    \node[text=blue!75!black,font=\scriptsize] at (0.28,1.87) {$b$};
  \end{scope}
  % -------- 2) y-constant plane --------
  \begin{scope}[shift={(8,0)}]
    \fill[orange!3,rounded corners=7pt] (-3.55,-3.15) rectangle (3.55,3.05);
    \draw[orange!20,rounded corners=7pt] (-3.55,-3.15) rectangle (3.55,3.05);
    \draw[->,gray!70] (-2.75,0)--(2.86,0) node[below,font=\scriptsize] {$x$};
    \draw[->,gray!70] (0,-2.48)--(0,2.55) node[left,font=\scriptsize] {$z$};
    \fill[orange!16] (0,0) circle [radius=1.84];
    \draw[orange!80!black,thick] (0,0) circle [radius=1.84];
    \fill (0,0) circle (1.2pt);
    \draw[orange!80!black,line width=1.2pt,->] (0,0)--(1.3,1.3);
    \node[font=\scriptsize,fill=orange!3] at (0,2.44) {$r=a\sqrt{1-\frac{y_0^2}{b^2}}$};
    \node[font=\scriptsize,text=orange!85!black] at (0,-2.25) {$y=y_0$：圆截面};
    \node[font=\scriptsize] at (0,-2.85) {$A(y_0)=\pi a^2(1-y_0^2/b^2)$};
  \end{scope}
  % -------- 3) x-constant plane --------
  \begin{scope}[shift={(16,0)}]
    \fill[teal!3,rounded corners=7pt] (-3.55,-3.15) rectangle (3.55,3.05);
    \draw[teal!17,rounded corners=7pt] (-3.55,-3.15) rectangle (3.55,3.05);
    \draw[->,gray!70] (-2.70,0)--(2.90,0) node[below,font=\scriptsize] {$y$};
    \draw[->,gray!70] (0,-2.44)--(0,2.60) node[left,font=\scriptsize] {$z$};
    \fill[teal!13] (0,0) ellipse [x radius=1.38,y radius=1.95];
    \draw[teal!75!black,thick] (0,0) ellipse [x radius=1.38,y radius=1.95];
    \draw[teal!75!black,->,line width=1.1pt] (0,0)--(1.38,0);
    \node[font=\scriptsize,text=teal!70!black,anchor=west] at (1.44,0.44) {$bq$};
    \draw[teal!75!black,->,line width=1.1pt] (0,0)--(0,1.95);
    \node[font=\scriptsize,text=teal!70!black,anchor=west] at (0.18,1.52) {$aq$};
    \node[font=\scriptsize,text=teal!75!black] at (0,-2.25) {$x=x_0$：椭圆截面};
    \node[font=\scriptsize] at (0,-2.85) {$q=\sqrt{1-x_0^2/a^2}$，$A=\pi abq^2$};
  \end{scope}
\end{tikzpicture}
}
    \end{center}
    \medskip
    \textbf{\textcolor{CExplanation}{新增视角二（球体伸缩：不用重新积分）}}

    先想象一个半径为 $b$ 的球：$x^2+y^2+z^2\le b^2$。保持 $y$ 方向不动，把\textbf{$x$、$z$ 两个方向同时}按 $a/b$ 倍拉伸（当 $a<b$ 时就是压缩）。球的三条半轴 $b,b,b$ 变成 $a,b,a$，恰好得到绕 $y$ 轴的旋转椭球。
    \[
       V_y=\underbrace{\frac43\pi b^3}_{\text{原球体积}}
       \underbrace{\frac ab\cdot\frac ab\cdot 1}_{\text{三方向伸缩倍数}}
       =\boxed{\frac43\pi a^2b}.
    \]
    \textbf{快记：} 球体积中的三个半径 $R\cdot R\cdot R$，对应换成椭球的三个半轴 $a\cdot a\cdot b$；统一系数 $\frac43\pi$ 不变。

\end{explanationbox}
